
\documentclass[12pt]{article}
\usepackage[margin=1in]{geometry}
\usepackage{setspace}
\usepackage{amsmath, amssymb}
\usepackage{natbib}
\usepackage{booktabs}
\usepackage{hyperref}

\title{Political Polarization and the Interpretation of Auditor Signals}
\author{Neil}
\date{\today}

\begin{document}
\maketitle
\onehalfspacing

\begin{abstract}
We examine whether political polarization affects how investors interpret auditor-related information. We propose that polarization increases dispersion in prior beliefs and reduces consensus about the credibility of auditors, leading to heterogeneous posterior beliefs following identical audit signals. Using auditor change events, we show that polarization amplifies trading volume, absolute returns, and disagreement. These findings contribute to auditing, capital markets, and political economy literatures.
\end{abstract}

\section{Introduction}

A central premise of auditing is that auditor reports enhance the credibility of financial information and facilitate efficient capital allocation. The accounting literature largely assumes that investors process audit signals in a relatively homogeneous manner. However, a growing body of research in finance and economics suggests that political polarization shapes beliefs, trust in institutions, and information processing.

This paper studies whether political polarization affects the interpretation of auditor signals. We argue that polarization generates heterogeneity in prior beliefs and trust in auditors, leading to differential interpretation of identical audit disclosures. As a result, polarization affects both price reactions and disagreement among investors.

\section{Literature Review}

\subsection{Audit and Capital Markets}

Early research demonstrates that audit opinions convey information to investors \citep{firth1978,dodd1984}. Subsequent work examines audit quality and market outcomes \citep{francis2004,defond2014}. More recent studies on expanded audit reporting provide mixed evidence on informativeness.

\subsection{Auditor Changes}

Auditor changes provide a salient signal of changes in audit quality or reporting risk \citep{johnson1990,defond1998,landsman2009}.

\subsection{Political Polarization}

Recent literature documents that political polarization affects financial markets, including trading behavior and asset prices \citep{kempf2024}.

\subsection{Disagreement and Information Processing}

Theoretical models show that heterogeneous beliefs lead to trading and volatility \citep{kandel1995,harris1993,hong2007}.

\section{Theory}

Investors observe a signal $s$ about firm quality $\theta$. Investors differ in priors $\pi_i$.

\begin{equation}
P_i(\theta|s) = \frac{P(s|\theta)\pi_i}{P(s)}
\end{equation}

Polarization increases dispersion in $\pi_i$, leading to greater disagreement.

Prices aggregate beliefs:

\begin{equation}
P = \sum_i w_i E_i[\theta|s]
\end{equation}

\textbf{Predictions:}
\begin{itemize}
\item Polarization increases absolute returns
\item Polarization increases trading volume
\item Effects stronger for ambiguous signals
\end{itemize}

\section{Empirical Design}

\subsection{Data}

We use CRSP, Compustat, and 13F data from WRDS, combined with SEC filings and election data.

\subsection{Variables}

Dependent variables:
\begin{itemize}
\item CAR (-1,+1)
\item Absolute CAR
\item Abnormal volume
\end{itemize}

Independent variables:
\begin{itemize}
\item Auditor change indicator
\item Polarization (county vote share dispersion, investor-based dispersion)
\end{itemize}

\subsection{Specification}

\begin{equation}
CAR_{i,t} = \alpha + \beta_1 Event_{i,t} + \beta_2 Polarization_{i,t} + \beta_3 Event \times Polarization + \gamma X + FE + \varepsilon
\end{equation}

\section{Identification Strategy}

We exploit cross-sectional variation in polarization and event-level variation. Firm and time fixed effects control for unobserved heterogeneity.

\section{Results Roadmap}

We estimate:
\begin{enumerate}
\item Baseline reactions
\item Heterogeneity by event type
\item Mechanism tests using investor composition
\item Robustness tests
\end{enumerate}

\section{Conclusion}

Political polarization affects the interpretation of audit signals, with implications for market efficiency and audit regulation.

\bibliographystyle{apalike}
\bibliography{final_refs}

\end{document}
